\documentclass[11pt]{article}
\usepackage[utf8]{inputenc}
\usepackage[T1]{fontenc}
\usepackage[a4paper,margin=2.5cm]{geometry}
\usepackage{amsmath,amssymb,graphicx,booktabs}
\usepackage{mathptmx}

% Rutas de figuras relativas a TP2/informe/, apuntando a TP2/data/plots/
\graphicspath{{../data/plots/}}

% Notacion (docs/GuiaInformes.md): vectores en negrita sin italica; los
% escalares quedan en italica por defecto de LaTeX en modo matematico.
\newcommand{\vect}[1]{\mathbf{#1}}

\title{Simulaci\'on de Bandadas: Modelo de Vicsek y Modelo de Votante}
\author{Grupo 5 -- Simulaci\'on de Sistemas}
\date{}

\begin{document}
\maketitle

\section{Introduccion}

Las bandadas de agentes autopropulsados son sistemas de muchas partículas que se
mueven a velocidad constante y ajustan su dirección en función de las direcciones
de sus vecinos cercanos, dando lugar a fenómenos colectivos de orden emergente sin
necesidad de un líder ni de reglas globales. Vicsek et al.\ [1]
propusieron un modelo mínimo off-lattice para este tipo de sistemas, en el cual
cada partícula promedia la dirección de las partículas dentro de un radio de
interacción $r_c$ y adopta esa dirección promedio perturbada por un ruido angular
$\eta$, mostrando que el sistema exhibe una transición de fase orden-desorden a
medida que $\eta$ aumenta.

Este trabajo implementa y caracteriza dicho modelo estándar de Vicsek, y lo compara
contra una regla de interacción alternativa conocida como modelo de votante
[2], en la cual cada partícula no promedia sino que copia
directamente la dirección de un único vecino elegido al azar (más el mismo ruido
$\eta$). El sistema se simula en una caja cuadrada de lado $L=10$ con condiciones
periódicas de contorno, para tres densidades $\rho = 2, 4, 8$, reutilizando el
Cell Index Method (CIM) desarrollado en el Trabajo Práctico 1 para la búsqueda
eficiente de vecinos dentro del radio $r_c$.

El objetivo del trabajo es estudiar el comportamiento del sistema en función del
parámetro de ruido $\eta$ para ambas reglas de interacción y las tres densidades
propuestas, caracterizando la transición orden-desorden a través de la
polarización $v_a$ (parámetro de orden global de las velocidades) y de la fracción
del cluster gigante $S$ (fracción de partículas que pertenecen a la componente
conexa más grande de la red de vecinos), comparando ambos modelos entre sí y
verificando además que el costo computacional del CIM reutilizado escala de forma
consistente con lo ya observado en el Trabajo Práctico 1.

\section{Modelo}

Cada partícula $i$ está descripta por su posición $\vect{r}_i(t)$ y su ángulo de
orientación $\theta_i(t)$, y se mueve a rapidez constante $v_0$. La posición se
actualiza de forma sincrónica (todas las partículas a la vez, a partir del estado
del paso anterior) según

\begin{equation}
  \vect{r}_i(t+\Delta t) = \vect{r}_i(t) + v_0 \left(\cos\theta_i(t+\Delta t),\
  \sin\theta_i(t+\Delta t)\right) \Delta t,
  \label{eq:position}
\end{equation}

con envoltura periódica de las coordenadas dentro de la caja de lado $L$ (condición
de contorno periódica). El vecindario $\mathcal{N}_i$ de la partícula $i$ es
autoinclusivo (contiene a $i$ y a toda partícula $j$ con $|\vect{r}_i - \vect{r}_j| < r_c$,
calculado en cada paso por el CIM) y es el mismo conjunto que usan ambas reglas de
interacción, difiriendo únicamente en cómo combinan las orientaciones de
$\mathcal{N}_i$.

En el \emph{modelo estándar de Vicsek} [1], la nueva orientación es el promedio
circular de las orientaciones de $\mathcal{N}_i$ más ruido angular:

\begin{equation}
  \theta_i(t+\Delta t) = \operatorname{atan2}\!\left(\sum_{j \in \mathcal{N}_i} \sin\theta_j(t),\;
  \sum_{j \in \mathcal{N}_i} \cos\theta_j(t)\right) + \Delta\theta_i,
  \label{eq:vicsek}
\end{equation}

donde la suma de senos y cosenos evita la discontinuidad de promediar ángulos
directamente. En el \emph{modelo de votante} [2], en cambio, la partícula no
promedia: elige un único vecino $j$ al azar dentro de $\mathcal{N}_i$ (con
distribución uniforme) y copia directamente su orientación:

\begin{equation}
  \theta_i(t+\Delta t) = \theta_j(t) + \Delta\theta_i,\qquad j \sim \mathcal{U}(\mathcal{N}_i).
  \label{eq:voter}
\end{equation}

En ambos modelos el ruido angular $\Delta\theta_i$ se agrega de la misma manera,
como una variable aleatoria uniforme $\Delta\theta_i \sim \mathcal{U}(-\eta/2,\ \eta/2)$,
independiente entre partículas y entre pasos, siendo $\eta$ el único parámetro que
controla el desorden del sistema. Los parámetros del sistema son la densidad
$\rho = N/L^2$, el radio de interacción $r_c$, la rapidez $v_0$ y el paso temporal
$\Delta t$, todos en unidades de simulación adimensionales (no se asocian a
unidades físicas explícitas, siguiendo la convención del enunciado y de [1]).

\section{Implementacion}

El motor de simulación (\texttt{TP2/src/}) reutiliza el Cell Index Method (CIM) del
Trabajo Práctico 1, adaptado a una nueva partícula \texttt{VicsekParticle} que
agrega orientación además de posición (\texttt{struct VicsekParticle \{ int id;
double x, y, theta; \}}, en \texttt{particle.h}). A diferencia del \texttt{computeCIM}
de TP1, que reconstruía la grilla desde cero en cada llamada, el nuevo tipo
\texttt{Grid} (\texttt{grid.h}) es un objeto persistente: sus buffers
\texttt{cells} y \texttt{neighbors\_} se reutilizan entre llamadas a
\texttt{rebuild()} en lugar de reasignarse en cada paso, eliminando la fricción de
reuso señalada tras TP1.

La clase \texttt{Simulation} (\texttt{engine/simulation.h}) encapsula el estado
completo de una corrida: el vector de partículas, la grilla persistente, un buffer
auxiliar \texttt{thetaNew\_} y el generador de números aleatorios. Su método
\texttt{step()} realiza una actualización sincrónica de dos buffers: primero
calcula la nueva orientación de \emph{todas} las partículas a partir de la
foto instantánea previa al paso (nunca mezclando orientaciones ya actualizadas de
otras partículas dentro del mismo paso), y recién después aplica esas orientaciones
para mover las posiciones, evitando sesgos por orden de iteración. La regla de
interacción es seleccionable en tiempo de construcción mediante un
\texttt{enum class Model \{ Vicsek, Voter \}}, que despacha a las funciones libres
\texttt{circularMeanHeading} o \texttt{voterHeading} (Ecs.~\ref{eq:vicsek} y
\ref{eq:voter} respectivamente); el ruido angular se agrega de forma idéntica en
ambos casos mediante una única función compartida \texttt{addAngularNoise}.

Los observables escalares se calculan directamente sobre la grilla del motor, sin
recomputar vecinos por fuera: \texttt{polarization()} calcula $v_a$ a partir del
vector de partículas, y \texttt{giantComponentFraction()} calcula $S$ recorriendo
las componentes conexas de \texttt{NeighborList} (la misma lista de adyacencia que
usa el CIM), ambas en \texttt{observables.h}. La salida se escribe en texto plano
(posición y velocidad real \texttt{vx, vy} por partícula y por paso, o un log
escalar reducido \texttt{(t, va, S)} para corridas de barrido), desacoplando la
simulación del módulo de animación en Python, tal como exige el enunciado. El
nuevo binario \texttt{tp2}/\texttt{tp2\_test} vive enteramente en \texttt{TP2/} y
reutiliza el CIM adaptándolo, sin modificar ningún archivo de \texttt{TP1/}.

\section{Simulaciones}

Todas las corridas del barrido paramétrico usan la geometría fija que impone el
enunciado: caja cuadrada de lado $L=10$ con condiciones periódicas de contorno. Se
estudian tres densidades $\rho \in \{2, 4, 8\}$, cada una con ambos modelos de
interacción (Vicsek y votante), en función del ruido angular $\eta$.

La grilla de valores de $\eta$ combina una grilla gruesa global de 9 puntos
equiespaciados en $[0, 2\pi]$ con una grilla fina de 8 puntos adicionales,
concentrada dentro del intervalo (\emph{bracket}) donde ocurre la transición
orden-desorden. Ese intervalo se ubica de forma independiente para cada
combinación (modelo, $\rho$) mediante un mini-barrido exploratorio previo (2
semillas y 500 pasos por punto, sobre la grilla gruesa), detectando el primer
cruce de la polarización media $v_a$ por debajo del umbral $0.5$. Esta ubicación
dinámica evita concentrar cómputo en una grilla fija que podría no coincidir con
la transición real de cada combinación.

Para cada punto $(\text{modelo}, \rho, \eta)$ de la grilla final se corren al menos
$K=5$ semillas independientes, con semilla determinística derivada de
$\mathrm{sha256}(\text{modelo}\,|\,\rho\,|\,\eta\,|\,\text{índice de repetición})$
truncado a 64 bits (nunca sembrada por reloj), garantizando reproducibilidad
exacta y decorrelación entre puntos de $\eta$ cercanos. Cada corrida completa del
barrido consta de 2000 pasos de simulación. El criterio de estado estacionario
descarta el primer 50\% de los pasos como transitorio, y esa misma ventana
temporal se usa para promediar tanto $v_a$ como $S$, de modo que ambos observables
se miden bajo el mismo criterio de corte.

\end{document}
